\section{Prompt pruning: как не задавить персонажа запретами}

Prompt pruning — уже упомянутый в 2.7 и 9.8 метод: вместо добавления нового правила на каждый нежелательный случай, находить архетипическую причину, из которой нежелательное поведение становится структурно невозможным, и одно предложение заменяет десяток частных запретов.

Механизм: если Рин где-то ведёт себя как servant, неправильный fix — добавить "не будь servant" как отдельное правило; правильный fix — усилить или прояснить формулировку её agency-позиции (5.6) так, чтобы servant-поведение противоречило уже написанному, а не отдельно запрещалось. Первое добавляет вес промпту без изменения глубинной причины; второе решает проблему на уровне источника.

Практический тест на необходимость pruning: если промпт растёт быстрее, чем растёт число реально предотвращённых проблем, это сигнал, что накапливаются симптоматические правила вместо причинных.
